% Theory for Project 02 — Paths, query parameters and request schemas.
% Build (from this folder):  tectonic theory.tex   ->  theory.pdf
\documentclass[11pt,a4paper]{article}
\usepackage[T1]{fontenc}
\usepackage{lmodern}
\usepackage[margin=22mm]{geometry}
\usepackage{booktabs}
\usepackage{tabularx}
\usepackage{array}
\usepackage{enumitem}
\usepackage[table]{xcolor}
\usepackage{fancyvrb}
\usepackage{titlesec}
\usepackage[colorlinks=true,linkcolor=black,urlcolor=accent]{hyperref}
\usepackage{xurl}

\definecolor{accent}{HTML}{3367D6}
\definecolor{boxbg}{HTML}{F4F6FA}
\titleformat{\section}{\large\bfseries\color{accent}}{\thesection}{0.6em}{}
\titleformat{\subsection}{\normalsize\bfseries}{\thesubsection}{0.6em}{}
\setlength{\parindent}{0pt}
\setlength{\parskip}{0.55em}
\setlist{itemsep=0.15em, topsep=0.3em}
\renewcommand{\arraystretch}{1.25}
\newcolumntype{L}{>{\raggedright\arraybackslash}X}
\newcommand{\code}[1]{\texttt{#1}}
\newcommand{\ra}{$\rightarrow$}

\newcommand{\keypoint}[1]{%
  \par\smallskip\noindent\colorbox{boxbg}{\parbox{\dimexpr\linewidth-2\fboxsep}{\small #1}}\par\smallskip}

\title{\textbf{Project 02 --- Paths, query parameters and request schemas}\\[0.3em]\large Theory, answers, design decision, scope and references}
\author{APIs: 24 Build-It-Yourself Projects}
\date{5 November 2026 \quad\textbullet\quad FastAPI 0.162.0, Pydantic 2.14.2 (core 2.50.2), Starlette 1.10.0, uvicorn 0.54.0, Python 3.13.18}

\begin{document}
\maketitle

\keypoint{\textbf{How to read this.} Every claim about \emph{our} API was checked
against the running code (TestClient, or a live uvicorn server called with curl on
2026-11-05; raw output in \code{evidence/live-2026-11-05.txt}). Claims about
libraries and standards link to the source in Section~\ref{sec:refs}. Where
something was measured rather than read, the text says so.}

\setcounter{tocdepth}{1}
\tableofcontents

% =====================================================================
\section{The boundary: types are a promise to clients}

An API's input types are not internal details. As soon as one client sends
\code{"quantity": "3"} and gets a 200, that spelling is part of the contract,
whatever the documentation says. This project therefore validates \emph{everything}
at the boundary, before the handler runs, and writes down every case where it still
coerces.

FastAPI decides where each value comes from by its declaration:

\begin{tabularx}{\linewidth}{@{}l L L@{}}
\toprule
\textbf{Declaration} & \textbf{Source} & \textbf{In our API} \\
\midrule
name inside \code{\{...\}} in the path & path segment & \code{product\_id: UUID} \\
plain type (\code{str}, \code{int}) not in the path & query string & none in P02 (P01's \code{name}) \\
Pydantic model & JSON body & \code{request: QuoteRequest} \\
\bottomrule
\end{tabularx}

The first element of an error's location \code{loc} follows the same split:
\code{path}, \code{query} or \code{body}, followed by the field name.

% =====================================================================
\section{Missing, null, invalid and zero are four different things}

The workbook's central claim. Measured on the live server for \code{quantity}:

\begin{tabularx}{\linewidth}{@{}l l l L@{}}
\toprule
\textbf{Body} & \textbf{Status} & \textbf{\code{type}} & \textbf{Meaning} \\
\midrule
\code{\{"unit\_price":"0.10"\}} & 422 & \code{missing} & the client did not say anything \\
\code{"quantity": null} & 422 & \code{int\_type} & the client said ``no value'' \\
\code{"quantity": 2.5} & 422 & \code{int\_type} & a value of the wrong type \\
\code{"quantity": 0} & 422 & \code{greater\_than\_equal} & right type, outside the range \\
\code{"quantity": 3} & 200 & --- & valid \\
\bottomrule
\end{tabularx}

\textbf{Required vs optional (R01).} In Pydantic a field is required exactly when it
has no default. \code{quantity: int = Field(ge=1)} is still required;
\code{Field(1, ge=1)} would silently make 1 the default, and an absent quantity would
be accepted.

\textbf{Null vs absent (R02).} An optional field written as \code{x: int | None = None}
would merge the two cases again, because both end up as \code{None}. When the
difference matters (for example in a PATCH, ``leave unchanged'' vs ``clear''), Pydantic
records which fields were actually sent in \code{model.model\_fields\_set}. P02 has no
optional fields, so both cases are rejected, with different error types.

% =====================================================================
\section{Money is not a float}

A binary float cannot store most decimal fractions exactly. Measured in Python 3.13:

\begin{Verbatim}[fontsize=\small, frame=single, framesep=2mm]
>>> 0.1 * 3
0.30000000000000004
>>> Decimal(0.1)        # the float's real value
0.1000000000000000055511151231257827021181583404541015625
>>> Decimal("0.10") * 3
Decimal('0.30')
\end{Verbatim}

JSON itself has only one number type, and RFC~8259 \S6 warns that implementations
commonly read numbers as IEEE~754 doubles. So the price travels as a \emph{string}
(\code{"0.10"}), and the server builds a \code{Decimal} from the exact digits.

\subsection{Rounding is a policy, not a default}

\code{quantize(Decimal("0.01"), rounding=...)} rounds to exactly two places. Python
offers eight rounding modes, and its default is \emph{half even}:

\begin{tabularx}{\linewidth}{@{}l l l L@{}}
\toprule
\textbf{Value} & \textbf{\code{ROUND\_HALF\_UP}} & \textbf{\code{ROUND\_HALF\_EVEN}} & \textbf{Note} \\
\midrule
0.125 & 0.13 & 0.12 & 2 is even, so half even goes down \\
0.135 & 0.14 & 0.14 & 4 is even, so half even goes up \\
\code{round(2.675, 2)} & \multicolumn{2}{l}{2.67 (a float, really 2.67499\ldots)} & why we never round floats \\
\bottomrule
\end{tabularx}

We chose half up because it matches what people expect on a receipt, and pinned it
with a test (\code{0.125} \ra\ \code{"0.13"}). Half even reduces bias over many
roundings and is common in accounting. Either is defensible. Leaving the choice
implicit is not.

\subsection{Bounds protect the arithmetic}

A \code{Decimal} operation runs in a \emph{context} with a precision of 28
significant digits by default:

\begin{Verbatim}[fontsize=\small, frame=single, framesep=2mm]
>>> Decimal("1E+100").quantize(Decimal("0.01"))   # needs 103 digits
decimal.InvalidOperation: [<class 'decimal.InvalidOperation'>]
\end{Verbatim}
 Before R04 our API turned
that into a \textbf{500} (found by probing, see EVIDENCE.md). With
\code{max\_digits=12, decimal\_places=4} the largest total is
$99\,999\,999.9999 \times 1000 \approx 10^{11}$, about 15 digits at most, so
\code{quantize} can no longer fail.

\keypoint{A number with no upper bound has not been validated. Bounds are needed for
the arithmetic, not only for the business rule.}

% =====================================================================
\section{Coercion: lax vs strict}

Pydantic validates in \emph{lax} mode by default: it converts values when the
conversion seems safe. For an \code{int} field, \code{"3"} and \code{3.0} become
\code{3}, and \code{2.5} is rejected (\code{int\_from\_float}). In \emph{strict} mode,
only a real integer is accepted. Our choices, all observed live:

\begin{tabularx}{\linewidth}{@{}l l L@{}}
\toprule
\textbf{Input} & \textbf{Result} & \textbf{Reason} \\
\midrule
\code{quantity: "3"}, \code{3.0} & 422 & \code{strict=True}: a count is an integer, nothing else \\
\code{unit\_price: "0.10"} & 200 & the intended form; must stay lax to parse a string \\
\code{unit\_price: 0.1} (number) & 200 & lax; the client has already lost precision \\
\code{"1E1"}, \code{" 0.10 "}, \code{"+0.10"}, \code{"1\_0"} & 200 & Python's \code{Decimal()} grammar accepts all of them \\
Arabic-Indic digits (U+0661 U+0662) & 200, total 12.00 & \code{Decimal()} accepts any Unicode decimal digit \\
\code{"-0"} & 200, \code{"0.00"} & $-0 \ge 0$ is true; the sign used to survive (\code{"-0.00"}), now normalized \\
\code{true}, \code{"0x10"} & 422 & not a decimal \\
\bottomrule
\end{tabularx}

\textbf{The trade-off.} RFC~9413 (\emph{Maintaining Robust Protocols}) argues against
the old ``be liberal in what you accept'' rule. Every tolerated variant becomes
something clients depend on, and you can never remove it without breaking them.
Strictness has the opposite cost: clients that send \code{"3"} break \emph{today}. We
chose strict where coercion could change meaning (a count) and documented the lax
cases we kept.

One side effect we measured: \code{strict=True} changed the error for \code{2.5} from
\code{int\_from\_float} to \code{int\_type}. Tests that assert on the library's
wording break on such changes. Ours assert on status and \code{loc}.

% =====================================================================
\section{Extra fields and mass assignment}

Pydantic's \code{extra} setting has three values: \code{"ignore"} (the default: unknown
keys are dropped silently), \code{"allow"} (they are kept on the model) and
\code{"forbid"} (422 \code{extra\_forbidden}).

\begin{itemize}
  \item \textbf{Ignore hides client bugs.} Our M1 typo \code{units\_price} produced
  only ``\code{unit\_price} missing''. With \code{forbid}, the client gets both
  ``\code{unit\_price} missing'' and ``\code{units\_price} not permitted'', and the typo
  is obvious.
  \item \textbf{Ignore misleads clients.} A client sending \code{"discount": "0.5"}
  receives a full-price quote with a 200, and may believe the discount was applied.
  \item \textbf{Allow invites mass assignment.} If extra keys are kept and copied into
  storage, a client can set fields it should never control (\code{is\_admin},
  \code{price}). OWASP lists this under API3:2023, \emph{Broken Object Property Level
  Authorization}.
\end{itemize}

\textbf{Measured limit:} \code{extra="forbid"} lives on the body model, so unknown
\emph{query} parameters are still ignored (\code{?admin=true} \ra\ 200). FastAPI can
also forbid extra query parameters by declaring them as a Pydantic model with
\code{Query()}. We did not do this in P02.

% =====================================================================
\section{Path parameters: 422 vs 404}

\code{product\_id: UUID} makes the path segment go through Python's \code{UUID()}
parser before the handler runs.

\begin{itemize}
  \item \code{/products/not-a-uuid} \ra\ \textbf{422} \code{uuid\_parsing} at
  \code{["path","product\_id"]}: this can \emph{never} identify a product. It's a client error.
  \item \code{/products/00000000-0000-4000-8000-000000000000} \ra\ \textbf{404}: well-formed,
  but nothing is stored under it.
\end{itemize}

RFC~9110 defines 404 as ``the origin server did not find a current representation for
the target resource'' (\S15.5.5) and 422 as ``unable to process the contained
instructions'' (\S15.5.21). Returning 404 for garbage would hide client bugs.
Returning 422 for a missing product would blame the client for a correct request.

A lesson from testing: \code{test\_unknown\_but\_valid\_product\_id\_is\_404} passed
\emph{before the route existed}, because every unknown path is a 404. A test is only
meaningful next to one that proves the route works (the 200 test).

\textbf{Accepted spellings} (measured): uppercase, no hyphens, \code{\{braces\}} and
\code{urn:uuid:} all reach the same product. The response is always the canonical
lowercase hyphenated form from RFC~9562. Because the dict keys are \code{UUID} objects,
not strings, the lookup is case-independent. P15 (caching) will need to decide
whether these are one URL or many.

% =====================================================================
\section{Internal vs public models}

R06 separates two schemas for one resource:

\begin{Verbatim}[fontsize=\small, frame=single, framesep=2mm]
class ProductRecord(BaseModel):     # internal: what we store
    id: UUID; name: str; supplier_password: str

class ProductResponse(BaseModel):   # public: what clients may see
    id: UUID; name: str

@app.get("/products/{product_id}", response_model=ProductResponse)
def get_product(product_id: UUID):
    ...
    return product                  # the whole record, on purpose
\end{Verbatim}

FastAPI re-validates the return value into \code{response\_model}, which drops
undeclared fields. The handler returns the risky object deliberately, so the test
shows that the \emph{schema} is the safety net, not the handler author's memory.

\textbf{Intentional failure (measured).} Without \code{response\_model}, FastAPI
serialized the whole record: 2 tests went red and printed the canary password.
Restored, all green again. A precondition test proves the secret really is stored;
otherwise ``the secret is absent'' would pass trivially.

\textbf{What \code{response\_model} does not cover.} Headers (P01's leak experiment),
logs, and the 422 \code{input} echo. For a \code{missing} error, \code{input} is the
\emph{whole request body}, so anything sensitive in it is sent back. This is carried
to P08.

% =====================================================================
\section{The workbook's four questions}

\textbf{1. What information crosses the contract boundary?} In: a JSON object with
exactly \code{unit\_price} (decimal string) and \code{quantity} (integer), or a UUID in
the path. Out: \code{total} as a 2-place decimal string, or \code{id} + \code{name}.
Errors: 422 with \code{type} + \code{loc} per problem, or 404. Nothing else
(\code{supplier\_password}, stack traces) crosses.

\textbf{2. Which invariant would survive a framework change?} ``The handler only ever
sees values inside the contract'', and ``the total is \code{unit\_price $\times$ quantity}
rounded half up to cents, as a decimal string''. The tests are written as HTTP
requests and responses, so they would still apply to another framework. Only the
\code{type} strings are Pydantic-specific.

\textbf{3. Which failure could hide behind a successful status?} Several did:
\code{"quantity": "3"} \ra\ 200 (silent coercion); \code{"discount"} ignored \ra\ 200
at full price; \code{"-0"} \ra\ 200 \code{"-0.00"} (now fixed); \code{?admin=true} ignored
\ra\ 200; and the 404 test, which was green while the route didn't exist.

\textbf{4. What had to be measured rather than read?} The 500 from \code{"1e100"};
exactly which spellings \code{Decimal()} and \code{UUID()} accept (underscores and
Unicode digits are not obvious from the docs); that \code{strict} changes error types;
and that \code{response\_model} filters a returned model of a \emph{different} class.

% =====================================================================
\section{Design decision}

\textbf{Chosen: declarative Pydantic constraints.} \code{Field(ge, le, max\_digits,
decimal\_places, strict)} + \code{extra="forbid"} + a separate response model.
\emph{Why:} validation is visible next to the field, runs before the handler, and
gives structured errors (\code{type}, \code{loc}) for free. \emph{Cost:} we inherit
the library's lax parsing (\code{Decimal()} grammar), and the \code{type} strings are
framework-specific.

\textbf{Alternative: a strict string type for \code{unit\_price}}, converted to
\code{Decimal} afterwards:

\begin{Verbatim}[fontsize=\small, frame=single, framesep=2mm]
unit_price: str = Field(pattern=r"^\d{1,8}(\.\d{1,4})?$")
\end{Verbatim}
 This would reject numbers, exponents, spaces,
\code{+}, \code{\_} and non-ASCII digits all at once. \emph{Cost:} a second
parsing step, and existing clients sending \code{0.1} as a number would break.

\textbf{What would make me switch:} the first time a lax spelling causes a real
disagreement (for example a downstream system that rejects \code{"1E1"}), or before a
public release, while no client depends on the lax forms yet.

% =====================================================================
\section{Scope of the project}

\textbf{In scope (built and tested):}
\begin{itemize}
  \item Two routes: \code{POST /quotes} (JSON body) and \code{GET /products/\{product\_id\}} (path).
  \item R01 required fields, R02 null vs absent, R03 decimal arithmetic + half-up
  rounding + strict integers, R04 numeric and digit bounds, R05 unknown body fields
  rejected, R06 internal fields kept out of responses.
  \item T01--T06, two boundary pairs, three regressions (\code{1e100}, \code{-0}, the
  \code{units\_price} typo), one intentional failure: 23 tests.
  \item A live check of 36 requests against uvicorn, and documented coercion rules.
\end{itemize}

\textbf{Deliberately out of scope:}
\begin{itemize}
  \item Persistence: \code{PRODUCTS} is an in-memory dict (P05).
  \item Custom error format: we keep FastAPI's 422 shape (P08).
  \item Query-parameter models and unknown query parameters (noted as open).
  \item Optional request fields and PATCH semantics (P03).
  \item Authentication and per-user visibility of fields (P09, P10).
  \item Latency measurement (methodology in P01), concurrency, CI.
\end{itemize}

\textbf{Known open items:} lax \code{unit\_price}
spellings; unknown query parameters ignored; 422 echoes \code{input}; \code{HEAD}
\ra\ 405 (from P01).

% =====================================================================
\section{References}\label{sec:refs}

\subsection*{Standards}
\begin{itemize}
  \item RFC 9110, \emph{HTTP Semantics}: \S15.5.5 404 Not Found, \S15.5.21 422 Unprocessable Content.\\
    \url{https://www.rfc-editor.org/rfc/rfc9110}
  \item RFC 8259, \emph{The JSON Data Interchange Format}: \S6 Numbers (interoperability and IEEE 754).\\
    \url{https://www.rfc-editor.org/rfc/rfc8259}
  \item RFC 9562, \emph{Universally Unique IDentifiers (UUIDs)}.\\
    \url{https://www.rfc-editor.org/rfc/rfc9562}
  \item RFC 9413, \emph{Maintaining Robust Protocols} (on Postel's law and tolerance).\\
    \url{https://www.rfc-editor.org/rfc/rfc9413}
\end{itemize}

\subsection*{Python}
\begin{itemize}
  \item \code{decimal}: contexts, rounding modes and \code{quantize}.\\
    \url{https://docs.python.org/3/library/decimal.html}
  \item Floating-Point Arithmetic: Issues and Limitations.\\
    \url{https://docs.python.org/3/tutorial/floatingpoint.html}
  \item \code{uuid}: UUID objects according to RFC 9562.\\
    \url{https://docs.python.org/3/library/uuid.html}
\end{itemize}

\subsection*{Pydantic}
\begin{itemize}
  \item Strict mode.\quad \url{https://docs.pydantic.dev/latest/concepts/strict_mode/}
  \item Conversion table (what lax and strict accept).\quad \url{https://docs.pydantic.dev/latest/concepts/conversion_table/}
  \item Fields and numeric / decimal constraints.\quad \url{https://docs.pydantic.dev/latest/concepts/fields/}
  \item Model configuration, \code{extra}.\quad \url{https://docs.pydantic.dev/latest/api/config/}
  \item Validation error types (\code{missing}, \code{int\_type}, \code{extra\_forbidden}, \ldots).\quad \url{https://docs.pydantic.dev/latest/errors/validation_errors/}
\end{itemize}

\subsection*{FastAPI}
\begin{itemize}
  \item Request body.\quad \url{https://fastapi.tiangolo.com/tutorial/body/}
  \item Path parameters.\quad \url{https://fastapi.tiangolo.com/tutorial/path-params/}
  \item Query parameter models (forbidding extra query parameters).\quad \url{https://fastapi.tiangolo.com/tutorial/query-param-models/}
  \item Response model, return type.\quad \url{https://fastapi.tiangolo.com/tutorial/response-model/}
  \item Handling errors.\quad \url{https://fastapi.tiangolo.com/tutorial/handling-errors/}
\end{itemize}

\subsection*{Security}
\begin{itemize}
  \item OWASP API Security Top 10 (2023), API3: Broken Object Property Level Authorization.\\
    \url{https://owasp.org/API-Security/editions/2023/en/0xa3-broken-object-property-level-authorization/}
  \item OWASP Mass Assignment Cheat Sheet.\\
    \url{https://cheatsheetseries.owasp.org/cheatsheets/Mass_Assignment_Cheat_Sheet.html}
\end{itemize}

\end{document}
